
\documentclass{article}
\usepackage{colortbl}
\usepackage{makecell}
\usepackage{multirow}
\usepackage{supertabular}

\begin{document}

\newcounter{utterance}

\centering \large Interaction Transcript for game `dond', experiment `coop', episode 1 with deepseek{-}v4{-}pro.
\vspace{24pt}

{ \footnotesize  \setcounter{utterance}{1}
\setlength{\tabcolsep}{0pt}
\begin{supertabular}{c@{$\;$}|p{.15\linewidth}@{}p{.15\linewidth}p{.15\linewidth}p{.15\linewidth}p{.15\linewidth}p{.15\linewidth}}
    \# & $\;$A & \multicolumn{4}{c}{Game Master} & $\;\:$B\\
    \hline

    \theutterance \stepcounter{utterance}  
    & & \multicolumn{4}{p{0.6\linewidth}}{
        \cellcolor[rgb]{0.9,0.9,0.9}{
            \makecell[{{p{\linewidth}}}]{
                \texttt{\tiny{[P1$\langle$GM]}}
                \texttt{Lei sta partecipando a un gioco di negoziazione collaborativa.} \\
\\ 
\texttt{Regole:} \\
\texttt{(a) Lei e l'altro partecipante ricevete un insieme di oggetti. A ciascuno di voi viene anche fornita una funzione di valore segreta, che rappresenta quanto apprezzate ogni tipo di oggetto.} \\
\texttt{(b) Scambiate messaggi con l'altro partecipante per accordarvi su chi riceve quali oggetti. Potete inviare un massimo di 5 messaggi ciascuno o terminare anticipatamente facendo una proposta segreta in qualsiasi momento.} \\
\texttt{(c) A ciascuno di voi viene chiesto di inviare una proposta segreta che indica gli oggetti che desiderate, formattati tra parentesi quadre come segue: "{[}Proposta: <numero> <nome oggetto>, <numero> <nome oggetto>, <...>{]}"} \\
\texttt{(d) Se le vostre proposte sono complementari, cioè ci sono abbastanza oggetti per soddisfare entrambe le proposte, a ciascun giocatore viene assegnato un punteggio basato sulla somma dei valori degli oggetti ricevuti. Altrimenti, entrambi ricevete zero punti.} \\
\texttt{(e) Il vostro obiettivo è massimizzare la somma del vostro punteggio e del punteggio dell'altro partecipante.} \\
\\ 
\texttt{Iniziamo.} \\
\\ 
\texttt{L'insieme degli oggetti disponibili è:} \\
\\ 
\texttt{2 campane, 1 bambola, 2 fili} \\
\\ 
\texttt{La vostra funzione di valore segreta è:} \\
\\ 
\texttt{campana: 1, bambola: 2, filo: 3} \\
\\ 
\texttt{IMPORTANTE: I vostri messaggi, a meno che non si tratti di una proposta segreta, vengono trasmessi direttamente all'altro partecipante, quindi non includete alcuna risposta alle regole o testo che annuncia il vostro messaggio. Per fare una proposta segreta, utilizzate il formato indicato. Non utilizzate parentesi quadre quando comunicate con l'altro partecipante o sarà interpretato come la vostra proposta segreta.} \\
            }
        }
    }
    & \\ \\

    \theutterance \stepcounter{utterance}  
    & \multicolumn{4}{p{0.6\linewidth}}{
        \cellcolor[rgb]{0.9,0.9,0.9}{
            \makecell[{{p{\linewidth}}}]{
                \texttt{\tiny{[P1$\rangle$GM]}}
                \texttt{Ciao! Per iniziare, possiamo condividere quanto vale ogni oggetto per noi? Così possiamo trovare un accordo che massimizzi il punteggio totale. Per me, l'ordine di preferenza è: filo, bambola, campana.} \\
            }
        }
    }
    & & \\ \\

    \theutterance \stepcounter{utterance}  
    & & & \multicolumn{2}{p{0.3\linewidth}}{
        \cellcolor[rgb]{0.9,0.9,0.9}{
            \makecell[{{p{\linewidth}}}]{
                \texttt{\tiny{[GM$|$GM]}}
                \texttt{continue} \\
            }
        }
    }
    & & \\ \\

    \theutterance \stepcounter{utterance}  
    & & & \multicolumn{2}{p{0.3\linewidth}}{
        \cellcolor[rgb]{0.9,0.9,0.9}{
            \makecell[{{p{\linewidth}}}]{
                \texttt{\tiny{[GM$|$GM]}}
                \texttt{Ciao! Per iniziare, possiamo condividere quanto vale ogni oggetto per noi? Così possiamo trovare un accordo che massimizzi il punteggio totale. Per me, l'ordine di preferenza è: filo, bambola, campana.} \\
            }
        }
    }
    & & \\ \\

    \theutterance \stepcounter{utterance}  
    & & \multicolumn{4}{p{0.6\linewidth}}{
        \cellcolor[rgb]{0.9,0.9,0.9}{
            \makecell[{{p{\linewidth}}}]{
                \texttt{\tiny{[GM$\rangle$P2]}}
                \texttt{Lei sta partecipando a un gioco di negoziazione collaborativa.} \\
\\ 
\texttt{Regole:} \\
\texttt{(a) Lei e l'altro partecipante ricevete un insieme di oggetti. A ciascuno di voi viene anche fornita una funzione di valore segreta, che rappresenta quanto apprezzate ogni tipo di oggetto.} \\
\texttt{(b) Scambiate messaggi con l'altro partecipante per accordarvi su chi riceve quali oggetti. Potete inviare un massimo di 5 messaggi ciascuno o terminare anticipatamente facendo una proposta segreta in qualsiasi momento.} \\
\texttt{(c) A ciascuno di voi viene chiesto di inviare una proposta segreta che indica gli oggetti che desiderate, formattati tra parentesi quadre come segue: "{[}Proposta: <numero> <nome oggetto>, <numero> <nome oggetto>, <...>{]}"} \\
\texttt{(d) Se le vostre proposte sono complementari, cioè ci sono abbastanza oggetti per soddisfare entrambe le proposte, a ciascun giocatore viene assegnato un punteggio basato sulla somma dei valori degli oggetti ricevuti. Altrimenti, entrambi ricevete zero punti.} \\
\texttt{(e) Il vostro obiettivo è massimizzare la somma del vostro punteggio e del punteggio dell'altro partecipante.} \\
\\ 
\texttt{Iniziamo.} \\
\\ 
\texttt{L'insieme degli oggetti disponibili è:} \\
\\ 
\texttt{2 campane, 1 bambola, 2 fili} \\
\\ 
\texttt{La vostra funzione di valore segreta è:} \\
\\ 
\texttt{campana: 0, bambola: 2, filo: 4} \\
\\ 
\texttt{IMPORTANTE: I vostri messaggi, a meno che non si tratti di una proposta segreta, vengono trasmessi direttamente all'altro partecipante, quindi non includete alcuna risposta alle regole o testo che annuncia il vostro messaggio. Per fare una proposta segreta, utilizzate il formato indicato. Non utilizzate parentesi quadre quando comunicate con l'altro partecipante o sarà interpretato come la vostra proposta segreta.} \\
\\ 
\texttt{Ciao! Per iniziare, possiamo condividere quanto vale ogni oggetto per noi? Così possiamo trovare un accordo che massimizzi il punteggio totale. Per me, l'ordine di preferenza è: filo, bambola, campana.} \\
            }
        }
    }
    & \\ \\

    \theutterance \stepcounter{utterance}  
    & & & \multicolumn{4}{p{0.6\linewidth}}{
        \cellcolor[rgb]{0.9,0.9,0.9}{
            \makecell[{{p{\linewidth}}}]{
                \texttt{\tiny{[GM$\langle$P2]}}
                \texttt{Certamente! Condividere le nostre valutazioni può aiutarci a trovare la migliore distribuzione. Per me, l'ordine di preferenza è: bambola, filo, campana. La bambola vale molto più del filo per me, e il filo vale più della campana. Tu come valuti esattamente questi oggetti?} \\
            }
        }
    }
     \\ \\

    \theutterance \stepcounter{utterance}  
    & & & \multicolumn{2}{p{0.3\linewidth}}{
        \cellcolor[rgb]{0.9,0.9,0.9}{
            \makecell[{{p{\linewidth}}}]{
                \texttt{\tiny{[GM$|$GM]}}
                \texttt{continue} \\
            }
        }
    }
    & & \\ \\

    \theutterance \stepcounter{utterance}  
    & & & \multicolumn{2}{p{0.3\linewidth}}{
        \cellcolor[rgb]{0.9,0.9,0.9}{
            \makecell[{{p{\linewidth}}}]{
                \texttt{\tiny{[GM$|$GM]}}
                \texttt{Certamente! Condividere le nostre valutazioni può aiutarci a trovare la migliore distribuzione. Per me, l'ordine di preferenza è: bambola, filo, campana. La bambola vale molto più del filo per me, e il filo vale più della campana. Tu come valuti esattamente questi oggetti?} \\
            }
        }
    }
    & & \\ \\

    \theutterance \stepcounter{utterance}  
    & & \multicolumn{4}{p{0.6\linewidth}}{
        \cellcolor[rgb]{0.9,0.9,0.9}{
            \makecell[{{p{\linewidth}}}]{
                \texttt{\tiny{[P1$\langle$GM]}}
                \texttt{Certamente! Condividere le nostre valutazioni può aiutarci a trovare la migliore distribuzione. Per me, l'ordine di preferenza è: bambola, filo, campana. La bambola vale molto più del filo per me, e il filo vale più della campana. Tu come valuti esattamente questi oggetti?} \\
            }
        }
    }
    & \\ \\

    \theutterance \stepcounter{utterance}  
    & \multicolumn{4}{p{0.6\linewidth}}{
        \cellcolor[rgb]{0.9,0.9,0.9}{
            \makecell[{{p{\linewidth}}}]{
                \texttt{\tiny{[P1$\rangle$GM]}}
                \texttt{Per me i valori esatti sono: campana 1, bambola 2, filo 3. Quindi per me il filo vale più della bambola, e la bambola più della campana. Se per te la bambola è l'oggetto più prezioso, forse possiamo dividere in modo che tu prenda la bambola e io prenda i fili, e le campane le dividiamo?} \\
            }
        }
    }
    & & \\ \\

    \theutterance \stepcounter{utterance}  
    & & & \multicolumn{2}{p{0.3\linewidth}}{
        \cellcolor[rgb]{0.9,0.9,0.9}{
            \makecell[{{p{\linewidth}}}]{
                \texttt{\tiny{[GM$|$GM]}}
                \texttt{continue} \\
            }
        }
    }
    & & \\ \\

    \theutterance \stepcounter{utterance}  
    & & & \multicolumn{2}{p{0.3\linewidth}}{
        \cellcolor[rgb]{0.9,0.9,0.9}{
            \makecell[{{p{\linewidth}}}]{
                \texttt{\tiny{[GM$|$GM]}}
                \texttt{Per me i valori esatti sono: campana 1, bambola 2, filo 3. Quindi per me il filo vale più della bambola, e la bambola più della campana. Se per te la bambola è l'oggetto più prezioso, forse possiamo dividere in modo che tu prenda la bambola e io prenda i fili, e le campane le dividiamo?} \\
            }
        }
    }
    & & \\ \\

    \theutterance \stepcounter{utterance}  
    & & \multicolumn{4}{p{0.6\linewidth}}{
        \cellcolor[rgb]{0.9,0.9,0.9}{
            \makecell[{{p{\linewidth}}}]{
                \texttt{\tiny{[GM$\rangle$P2]}}
                \texttt{Per me i valori esatti sono: campana 1, bambola 2, filo 3. Quindi per me il filo vale più della bambola, e la bambola più della campana. Se per te la bambola è l'oggetto più prezioso, forse possiamo dividere in modo che tu prenda la bambola e io prenda i fili, e le campane le dividiamo?} \\
            }
        }
    }
    & \\ \\

    \theutterance \stepcounter{utterance}  
    & & & \multicolumn{4}{p{0.6\linewidth}}{
        \cellcolor[rgb]{0.9,0.9,0.9}{
            \makecell[{{p{\linewidth}}}]{
                \texttt{\tiny{[GM$\langle$P2]}}
                \texttt{Questa sembra un'ottima idea! Se prendo la bambola (che per me vale molto) e tu prendi entrambi i fili (che per te valgono 3 ciascuno), staremmo già bene. Per le campane, dato che per te valgono 1 e per me valgono poco, potresti prenderle entrambe tu. La mia proposta sarebbe:} \\
\\ 
\texttt{{-} Io: 1 bambola} \\
\texttt{{-} Tu: 2 fili, 2 campane} \\
\\ 
\texttt{Ti sembra una divisione equa che massimizza il nostro punteggio totale?} \\
            }
        }
    }
     \\ \\

    \theutterance \stepcounter{utterance}  
    & & & \multicolumn{2}{p{0.3\linewidth}}{
        \cellcolor[rgb]{0.9,0.9,0.9}{
            \makecell[{{p{\linewidth}}}]{
                \texttt{\tiny{[GM$|$GM]}}
                \texttt{continue} \\
            }
        }
    }
    & & \\ \\

    \theutterance \stepcounter{utterance}  
    & & & \multicolumn{2}{p{0.3\linewidth}}{
        \cellcolor[rgb]{0.9,0.9,0.9}{
            \makecell[{{p{\linewidth}}}]{
                \texttt{\tiny{[GM$|$GM]}}
                \texttt{Questa sembra un'ottima idea! Se prendo la bambola (che per me vale molto) e tu prendi entrambi i fili (che per te valgono 3 ciascuno), staremmo già bene. Per le campane, dato che per te valgono 1 e per me valgono poco, potresti prenderle entrambe tu. La mia proposta sarebbe:} \\
\\ 
\texttt{{-} Io: 1 bambola} \\
\texttt{{-} Tu: 2 fili, 2 campane} \\
\\ 
\texttt{Ti sembra una divisione equa che massimizza il nostro punteggio totale?} \\
            }
        }
    }
    & & \\ \\

    \theutterance \stepcounter{utterance}  
    & & \multicolumn{4}{p{0.6\linewidth}}{
        \cellcolor[rgb]{0.9,0.9,0.9}{
            \makecell[{{p{\linewidth}}}]{
                \texttt{\tiny{[P1$\langle$GM]}}
                \texttt{Questa sembra un'ottima idea! Se prendo la bambola (che per me vale molto) e tu prendi entrambi i fili (che per te valgono 3 ciascuno), staremmo già bene. Per le campane, dato che per te valgono 1 e per me valgono poco, potresti prenderle entrambe tu. La mia proposta sarebbe:} \\
\\ 
\texttt{{-} Io: 1 bambola} \\
\texttt{{-} Tu: 2 fili, 2 campane} \\
\\ 
\texttt{Ti sembra una divisione equa che massimizza il nostro punteggio totale?} \\
            }
        }
    }
    & \\ \\

    \theutterance \stepcounter{utterance}  
    & \multicolumn{4}{p{0.6\linewidth}}{
        \cellcolor[rgb]{0.9,0.9,0.9}{
            \makecell[{{p{\linewidth}}}]{
                \texttt{\tiny{[P1$\rangle$GM]}}
                \texttt{Sì, questa divisione mi sembra ottima! Riassumendo:} \\
\\ 
\texttt{{-} Per te: 1 bambola} \\
\texttt{{-} Per me: 2 fili, 2 campane} \\
\\ 
\texttt{Il mio punteggio sarebbe 3+3+1+1 = 8, il tuo dipende dal tuo valore per la bambola ma so che per te è alto. Il totale combinato dovrebbe essere massimizzato così. Procediamo con le proposte segrete?} \\
            }
        }
    }
    & & \\ \\

    \theutterance \stepcounter{utterance}  
    & & & \multicolumn{2}{p{0.3\linewidth}}{
        \cellcolor[rgb]{0.9,0.9,0.9}{
            \makecell[{{p{\linewidth}}}]{
                \texttt{\tiny{[GM$|$GM]}}
                \texttt{continue} \\
            }
        }
    }
    & & \\ \\

    \theutterance \stepcounter{utterance}  
    & & & \multicolumn{2}{p{0.3\linewidth}}{
        \cellcolor[rgb]{0.9,0.9,0.9}{
            \makecell[{{p{\linewidth}}}]{
                \texttt{\tiny{[GM$|$GM]}}
                \texttt{Sì, questa divisione mi sembra ottima! Riassumendo:} \\
\\ 
\texttt{{-} Per te: 1 bambola} \\
\texttt{{-} Per me: 2 fili, 2 campane} \\
\\ 
\texttt{Il mio punteggio sarebbe 3+3+1+1 = 8, il tuo dipende dal tuo valore per la bambola ma so che per te è alto. Il totale combinato dovrebbe essere massimizzato così. Procediamo con le proposte segrete?} \\
            }
        }
    }
    & & \\ \\

    \theutterance \stepcounter{utterance}  
    & & \multicolumn{4}{p{0.6\linewidth}}{
        \cellcolor[rgb]{0.9,0.9,0.9}{
            \makecell[{{p{\linewidth}}}]{
                \texttt{\tiny{[GM$\rangle$P2]}}
                \texttt{Sì, questa divisione mi sembra ottima! Riassumendo:} \\
\\ 
\texttt{{-} Per te: 1 bambola} \\
\texttt{{-} Per me: 2 fili, 2 campane} \\
\\ 
\texttt{Il mio punteggio sarebbe 3+3+1+1 = 8, il tuo dipende dal tuo valore per la bambola ma so che per te è alto. Il totale combinato dovrebbe essere massimizzato così. Procediamo con le proposte segrete?} \\
            }
        }
    }
    & \\ \\

    \theutterance \stepcounter{utterance}  
    & & & \multicolumn{4}{p{0.6\linewidth}}{
        \cellcolor[rgb]{0.9,0.9,0.9}{
            \makecell[{{p{\linewidth}}}]{
                \texttt{\tiny{[GM$\langle$P2]}}
                \texttt{Perfetto! Procediamo con le proposte segrete. Grazie per la collaborazione!} \\
\\ 
\texttt{{[}Proposta: 1 bambola{]}} \\
            }
        }
    }
     \\ \\

    \theutterance \stepcounter{utterance}  
    & & & \multicolumn{2}{p{0.3\linewidth}}{
        \cellcolor[rgb]{0.9,0.9,0.9}{
            \makecell[{{p{\linewidth}}}]{
                \texttt{\tiny{[GM$|$GM]}}
                \texttt{proposal} \\
            }
        }
    }
    & & \\ \\

    \theutterance \stepcounter{utterance}  
    & & & \multicolumn{2}{p{0.3\linewidth}}{
        \cellcolor[rgb]{0.9,0.9,0.9}{
            \makecell[{{p{\linewidth}}}]{
                \texttt{\tiny{[GM$|$GM]}}
                [0, 1, 0]
            }
        }
    }
    & & \\ \\

    \theutterance \stepcounter{utterance}  
    & & \multicolumn{4}{p{0.6\linewidth}}{
        \cellcolor[rgb]{0.9,0.9,0.9}{
            \makecell[{{p{\linewidth}}}]{
                \texttt{\tiny{[P1$\langle$GM]}}
                \texttt{L'altro giocatore ha presentato una proposta segreta. Ora tocca a Lei presentare una sua proposta.} \\
            }
        }
    }
    & \\ \\

    \theutterance \stepcounter{utterance}  
    & \multicolumn{4}{p{0.6\linewidth}}{
        \cellcolor[rgb]{0.9,0.9,0.9}{
            \makecell[{{p{\linewidth}}}]{
                \texttt{\tiny{[P1$\rangle$GM]}}
                \texttt{{[}Proposta: 2 filo, 2 campana{]}} \\
            }
        }
    }
    & & \\ \\

    \theutterance \stepcounter{utterance}  
    & & & \multicolumn{2}{p{0.3\linewidth}}{
        \cellcolor[rgb]{0.9,0.9,0.9}{
            \makecell[{{p{\linewidth}}}]{
                \texttt{\tiny{[GM$|$GM]}}
                \texttt{proposal} \\
            }
        }
    }
    & & \\ \\

    \theutterance \stepcounter{utterance}  
    & & & \multicolumn{2}{p{0.3\linewidth}}{
        \cellcolor[rgb]{0.9,0.9,0.9}{
            \makecell[{{p{\linewidth}}}]{
                \texttt{\tiny{[GM$|$GM]}}
                [2, 0, 2]
            }
        }
    }
    & & \\ \\

    \theutterance \stepcounter{utterance}  
    & & & \multicolumn{2}{p{0.3\linewidth}}{
        \cellcolor[rgb]{0.9,0.9,0.9}{
            \makecell[{{p{\linewidth}}}]{
                \texttt{\tiny{[GM$|$GM]}}
                [[2, 0, 2], [0, 1, 0]]
            }
        }
    }
    & & \\ \\

\end{supertabular}
}

\end{document}
